\section{Layer L1: information and interface models}
\label{sec:layer1}

L1 is where industrial automation has invested most and where an agent finds the
richest \emph{type} information. In a few important cases it also holds the only
machine-readable statement of what may legally be \emph{done}.

\subsection{\opcua{} (IEC 62541)}

\opcua{} \cite{iec62541} is the dominant vendor-neutral information model of
factory and process automation. Its core is a typed, navigable \emph{address space}: nodes
(objects, variables, methods, types) connected by typed references
(\texttt{HasComponent}, \texttt{HasProperty}, \texttt{Organizes},
\texttt{HasTypeDefinition}). Object-oriented modelling, with type definitions
that carry inheritance and are instantiated as objects, lets a client discover
the structure of a machine it has never seen, at run time, without a schema
shipped out of band.
\emph{Companion specifications} standardise the model for a domain (robotics,
machine tools, machine vision, packaging via PackML, and several dozen others),
so that two robots from different vendors expose the same browse names for the
same concepts. Serialisation as \texttt{NodeSet2} XML makes the whole model a
transferable artefact, and PubSub extends the transport beyond client/server.

\paragraph{Agent affordances.} Four are decisive.
\emph{(i) Introspectable typing.} \texttt{DataType}, \texttt{ValueRank} and
\texttt{ArrayDimensions} let an agent emit correctly-shaped values.
\emph{(ii) Access control as data.} \texttt{AccessLevel} /
\texttt{UserAccessLevel} distinguish readable from writable nodes. This is
what makes fault family \textsc{f-acc} detectable at all, and it accounts for
eight of the caught cases in Experiment~E4.
\emph{(iii) Engineering ranges.} \texttt{EURange} and \texttt{InstrumentRange}
give constrained decoding real bounds; \texttt{EUInformation} gives a unit
display string.
\emph{(iv) Methods.} \opcua{} \texttt{Method} nodes with typed input and output
arguments are, structurally, function-calling tool schemas that already exist in
the plant. Projecting them into an agent's tool list is close to mechanical, and
is the single highest-leverage integration available today.

\paragraph{Limitations for agents.} \opcua{} units are an
\texttt{EUInformation} structure referencing UN/CEFACT codes; there is no
dimension algebra and no conversion, so \texttt{bar} and \texttt{psi} are two
opaque strings (U3~$=1$). Cross-machine topology is not modelled: the address
space is a containment tree per server, not a plant graph, so ``which units lose
air if the compressor trips'' is unanswerable within \opcua{} alone
(U4~$=2$). Constraints beyond range and type are prose in the companion
specification (U6~$=1$). And \texttt{NodeSet2} is extraordinarily verbose:
Experiment~E3 measures our small plant at over $23\times$ the token cost of the
equivalent tag list, so an agent must never be handed one directly.

\subsection{The Asset Administration Shell (IDTA)}

The AAS \cite{idta2023aas} is the Industrie~4.0 digital-twin metamodel,
standardised by the Industrial Digital Twin Association in a multi-part
specification (metamodel, API, IEC~61360 data specification, security, and the
AASX package format). An
asset has a shell; a shell has \emph{submodels}; a submodel contains
\emph{submodel elements}, which may be properties, collections, ranges, files,
relationship elements or operations. Every element may carry a
\texttt{semanticId} (typically an IRDI or IRI), which is what binds the AAS to
L0. Standardised \emph{submodel templates}, such as Digital Nameplate and
Technical Data \cite{idta2023nameplate} together with a growing library covering
documentation, carbon footprint and spare parts, make the same information
appear in the same shape across vendors. The familiar Type~1/2/3
distinction (a file, a served but passive shell, an actively communicating
shell) is a deployment maturity axis rather than a modelling one.

\paragraph{Agent affordances.} The AAS is the best-aligned industrial standard
for agent consumption, for three reasons. First, submodel templates give an
agent a \emph{stable expectation}: if the asset has a Nameplate submodel, the
manufacturer is at a known \texttt{idShort} with a known \texttt{semanticId}, so
extraction becomes lookup rather than inference. Second, the REST API
(Part~2) is a natural tool surface: \texttt{GET /submodels/\{id\}/\allowbreak
submodel-elements/\{idShortPath\}} is a function call with a typed return.
Third, \texttt{semanticId} plus the IEC~61360 data specification carries unit and
datatype down to the individual property. In our reference plant the AAS
projection is what condition~C2 renders.

\paragraph{Limitations for agents.} The AAS models an \emph{asset} and is weak
at modelling the \emph{plant}: cross-asset relations exist
(\texttt{RelationshipElement}, \texttt{AnnotatedRelationshipElement}) but are
neither standardised into a topology vocabulary nor transitively queryable.
There is no query language in the metamodel; the specification is currently
extending in this direction, but as of writing an agent must fetch and filter
client-side, which is exactly the pattern Experiment~E7 shows to be a
$17$--$\MaxQuerySaving\times$ token penalty. Behavioural semantics are limited to
\texttt{Operation} elements; there is no standard state model. And the serialised
environment is the most verbose artefact we measured ($29\times$ the tag list).

\subsection{AutomationML (IEC 62714)}

AutomationML \cite{iec62714} is the engineering-data exchange format: CAEX for
topology, with
libraries of system-unit classes, role classes and interface classes instantiated
into an instance hierarchy; COLLADA for geometry
and kinematics; PLCopen~XML for behaviour. Its \emph{role classes} are its
semantic contribution: an object is declared to play a role
(\texttt{Robot}, \texttt{Conveyor}) independently of the vendor class that
implements it.

\paragraph{Agent affordance.} AutomationML has the strongest \emph{as-engineered}
relational model of any L1 technology (U4~$=3$): interfaces and links express
plant topology explicitly, which is exactly the structure that Experiment~E2
shows to be decisive for impact-analysis tasks (TOP-01--05). Converters between
AutomationML and the AAS are mature, and the pragmatic pattern of lifting
engineering topology out of AutomationML and publishing it as AAS submodels or
RDF is one of the cheapest routes to a plant graph.

\paragraph{Limitation.} It is an engineering-time artefact. It does not carry
run-time values, states or events, so on its own it can support planning but not
situational awareness.

\subsection{MTP (VDI/VDE/NAMUR 2658)}

The Module Type Package \cite{vdi2658} describes a process module as a set of
\emph{services} with parameters, and a state machine per service, plus an HMI
description and communication mapping (typically \opcua). It exists because
modular process plants must be orchestrated by a Process Orchestration Layer
that has never seen the module before.

\paragraph{Agent affordance.} MTP is, from an agent's point of view, unusually
well designed: it is one of the very few industrial standards that describes a
plant component as a \emph{set of invocable services with legal state
transitions and typed parameters}. That is a tool schema with a precondition
model. An agent orchestrating MTP modules can plan without a hand-written
capability description, and can be prevented from issuing an illegal service
command by construction. It scores $3$ on U5 and $2$ on U6, a combination almost
nothing else achieves.

\paragraph{Limitation.} Adoption is concentrated in process industries, and the
service model does not extend to the discrete world.

\subsection{PackML / ISA-88 and ISA-95 / IEC 62264}

PackML (ANSI/ISA-TR88.00.02 \cite{isa88}, and the corresponding \opcua{}
companion model) standardises a machine's operating states, namely
\texttt{Idle}, \texttt{Execute}, \texttt{Held}, \texttt{Stopped},
\texttt{Aborted} and the acting states between them, together with the commands
that move between them. ISA-95 /
IEC~62264 \cite{iec62264} standardises the enterprise-to-control hierarchy
(enterprise, site, area, work centre, work unit) and the object models for
production schedules, definitions and performance, with B2MML as its XML
binding.

\paragraph{Agent affordance.} These two are small, cheap and disproportionately
valuable. PackML is the entire content of criterion U5 for packaging and
discrete lines: it turns ``can I start the filler?'' from a judgement into a
lookup, and it is what makes fault family \textsc{f-state} catchable. That
family accounts for seven cases in our corpus and is undetectable by every other
tier. PackML is also
\emph{token-cheap}: the complete state machine is a few hundred tokens, the best
value-per-token of any artefact we measured. ISA-95 supplies the mereology that
scopes every query an agent asks (``the Line~1 tags'', ``the utilities area
total''), and its absence is why fault family \textsc{f-scope} is invisible below
tier~G3.

\paragraph{Limitation.} Neither carries units, values or constraints. PackML in
particular is frequently \emph{implemented} without being \emph{published}: the
states exist in the PLC, but nothing exposes them as a model, which is the
practical reason so many deployed agents cannot use them.

\subsection{DTDL, W3C WoT Thing Description, NGSI-LD}

Three technologies converge from three directions on the same shape, a
JSON(-LD) description of a thing's properties, actions or commands, and events.

\emph{DTDL} (Digital Twins Definition Language, Azure Digital Twins)
\cite{dtdl} offers interfaces with telemetry, properties, commands, components
and typed relationships, and a graph query language over the twin graph. It is
the strongest of the three on relations plus query, and the weakest on openness.

\emph{W3C WoT Thing Description~1.1} \cite{w3c2023wot} is a W3C Recommendation
and, for our purposes, the most elegant: properties, actions and events each
carry a JSON Schema data model \emph{and} protocol bindings (\texttt{forms}),
while the JSON-LD context allows every affordance to be annotated with an
external ontology term. A Thing Description is very nearly a function-calling
manifest with semantic annotations attached, and is the closest existing
standard to what an agent tool registry should look like.

\emph{NGSI-LD} (ETSI GS~CIM~009 \cite{etsi2023ngsild}) models entities with
properties and relationships, where every attribute may itself carry metadata
(unit code, \texttt{observedAt} timestamp, \texttt{datasetId}) and the whole is
grounded in JSON-LD/RDF. Its provenance handling is the best in L1 (U7~$=3$), and its
\texttt{observedAt} per attribute is exactly what fault family
\textsc{f-stale} requires.

\paragraph{Assessment.} Together these show what a modern industrial standard
\emph{could} look like: typed affordances, external semantic annotation, native
JSON-LD, timestamps as first-class attribute metadata. Their weakness is
domain depth. None carries a plant taxonomy, an equipment state model or an
engineering-range convention, so they are complements to \opcua/AAS rather than
replacements.

\subsection{Sector-specific models: MTConnect, IEC 61850, Sparkplug B}

MTConnect \cite{mtconnect} (machine tools) offers a clean, read-oriented XML
model with \texttt{DataItem}s that carry a type, a subtype, units and native
units. It is one of the few standards that distinguishes reporting units from
native units, which is a small but real U3 contribution. IEC~61850
\cite{iec61850} (substation automation) has arguably the most rigorous naming
and data-model discipline in all of automation: logical nodes and standardised
data objects give near-perfect L0/L1 alignment inside its domain. Sparkplug~B
\cite{sparkplug} is included as the pragmatic counterpoint: it
adds a birth-certificate metadata payload to MQTT, giving a minimal but
genuinely useful type and unit hint at very low token cost. It is a
reminder that a modest, universally deployed semantic increment can be worth more
than a rich, rarely deployed one.

\subsection{Summary of L1}

\begin{table*}[t]
\centering
\caption{What each L1 technology uniquely contributes to an agent, and what it
cannot do. ``Unique'' means no other widely deployed L1 standard supplies it.}
\label{tab:l1-summary}
\footnotesize
\begin{tabular}{@{}lp{5.6cm}p{5.6cm}@{}}
\toprule
\textbf{Technology} & \textbf{Unique agent contribution} & \textbf{Principal gap} \\
\midrule
\opcua{} & Run-time introspectable types, access rights, EU ranges, callable
methods & No unit calculus, no plant-wide topology, extreme verbosity \\
AAS & Standardised submodel templates + \texttt{semanticId} down to the property;
REST tool surface & No query language, weak topology, no state model \\
AutomationML & As-engineered topology and role classes & Engineering-time only \\
MTP & Services with parameters \emph{and} legal state transitions & Process
industries only \\
PackML & The legal command set in the current state, at negligible token cost &
Nothing else \\
ISA-95 & The mereology that scopes every query & No values, units or constraints \\
WoT TD & Affordances with protocol bindings and external semantic annotation &
No domain depth \\
NGSI-LD & Per-attribute unit code and \texttt{observedAt} provenance & Weak
behaviour model \\
DTDL & Twin graph with a query language & Vendor-bound \\
\bottomrule
\end{tabular}
\end{table*}
